\section{Síntesis y ampliación}

\subsection{Tres principios fundamentales}

Anthropic resume en la conclusión del artículo los tres principios fundamentales para construir un agente:

\begin{importantbox}{Tres principios fundamentales para construir un agente}
\begin{enumerate}
    \item \textbf{Mantener el diseño sencillo} (Maintain simplicity): La arquitectura del agente debe ser lo más simple posible, incrementando la complejidad únicamente cuando se demuestre que es efectiva.
    \item \textbf{Priorizar la transparencia}: Exponer claramente los pasos de planificación del agente para que tanto usuarios como desarrolladores comprendan el proceso de decisión del agente.
    \item \textbf{Diseñar cuidadosamente la ACI}: A través de una documentación exhaustiva de las herramientas y pruebas, garantizar que el agente interactúe correctamente y eficientemente con el mundo exterior.
\end{enumerate}
\end{importantbox}

\subsection{Qué aprender de este artículo}

El valor de este artículo no reside únicamente en los patrones de diseño específicos, sino en la \textbf{filosofía de ingeniería} que transmite:

\begin{enumerate}
    \item \textbf{De lo simple a lo complejo}: No dejarse confundir por el concepto de «agente»; intentar primero con la solución más simple y luego escalar progresivamente.
    \item \textbf{Pensamiento basado en patrones}: Los cinco patrones de flujo de trabajo (workflows) y el patrón de agente ofrecen un marco de referencia práctico que ayuda a los desarrolladores a emparejar rápidamente la arquitectura adecuada al enfrentar nuevas tareas.
    \item \textbf{Las herramientas son la interfaz}: El límite superior de las capacidades del agente depende en gran medida de la calidad del diseño de las herramientas, y no tanto de el modelo mismo ni de la sofisticación del prompt.
    \item \textbf{Colaboración humano-máquina}: Incluso para un agente altamente autónomo, se debe diseñar un mecanismo de supervisión e intervención humana.
\end{enumerate}

\subsection{Lecturas adicionales}

\begin{itemize}
    \item Cookbook oficial de Anthropic: \url{https://github.com/anthropics/anthropic-cookbook} — Incluye ejemplos de implementación de código para diversos patrones agentic.
    \item Model Context Protocol (MCP): \url{https://modelcontextprotocol.io/} — El estándar unificado para el ecosistema de herramientas de agentes.
    \item SWE-bench: \url{https://www.swebench.com/} — Un benchmark estándar para evaluar las capacidades de los agentes de programación.
    \item Claude Agent SDK: \url{https://github.com/anthropics/claude-agent-sdk} — El kit de desarrollo oficial de Anthropic para agentes.
\end{itemize}

%% --- Fin del contenido --- %%

\end{document}
